\chapter{图件总表}
\label{ch:figures}

成稿的图分为机制图、数学图和数据图。机制图表达变量之间的依存，数学图表达证明或病态性的几何，数据图只使用第 \ref{ch:sim} 章日志中的列。日志未生成的图，在正文中保留坐标说明和空白图位，不绘制示意曲线充当结果。

全稿图位按本章清单实施。基础部分为机制图 28 幅、数学图 16 幅、数据图 24 幅、实施例 8 幅；S1 与 S2 的扩展面板另计 D25--D36。合计 88 幅独立图号。分面板用 (a)(b)(c)，不另计科学内容。

\section{机制图}

\begin{table}[H]
\centering
\caption{机制图 M1--M28}
\begin{tabularx}{\textwidth}{clX}
\toprule
编号 & 归属 & 图中必须出现的对象 \\
\midrule
M1 & 第 1 章 & 四个决策与 $\Delta\mathcal{I}$ 的对应 \\
M2 & 第 1 章 & 六步闭环及真值隔离的切口 \\
M3 & 第 3 章 & 五层与唯一执行器接口 \\
M4 & 第 3 章 & 四个时间尺度的中断关系 \\
M5 & 第 3 章 & 模块合同：只传递分布 \\
M6 & 第 4 章 & 固、液、汽三相与界面 $\Gamma$ \\
M7 & 第 4 章 & 能量方程中的对流、扩散、潜热、反应热 \\
M8 & 第 4 章 & 冷凝与蒸发对 $q_{\mathrm{lat}}$ 符号的切换 \\
M9 & 第 4 章 & 蛋白、淀粉、美拉德、热力当量四条进度 \\
M10 & 第 4 章 & 结壳增大 $R_{\mathrm{int}}$ 的反馈 \\
M11 & 第 4 章 & 全场 $\vect{X}$ 到决策状态 $\vect{z}$ 的映射 $\Phi$ \\
M12 & 第 4 章 & GAB 等温线在两类食材上的形状差异 \\
M13 & 第 4 章 & $k_\parallel$、$k_\perp$ 与脂肪夹层的热阻位置 \\
M14 & 第 4 章 & 收缩使厚度与渗透率同时变化 \\
M15 & 第 4 章 & 加热管、腔壁、食品之间的角系数 \\
M16 & 第 5 章 & 红外亮度中的发射与反射 \\
M17 & 第 5 章 & 冷点偏离几何中心的两种边界 \\
M18 & 第 6 章 & P4 的两段执行器时序 \\
M19 & 第 6 章 & 表面时间常数与冷点 Fourier 窗口 \\
M20 & 第 6 章 & 对偶目标中信息项随赤字开闭 \\
M21 & 第 7 章 & 父节点固定的结构图 \\
M22 & 第 7 章 & 路径 A、B、C 的分叉与汇合 \\
M23 & 第 7 章 & 预测、干预、反事实三个计算 \\
M24 & 第 8 章 & 物理项、有界残差、越界回退 \\
M25 & 第 8 章 & 类别先验的精度加权，含拒绝写回 \\
M26 & 第 9 章 & 候选执行量到安全集的投影 \\
M27 & 第 9 章 & 五级降级只向更少依赖估计的方向 \\
M28 & 第 9 章 & 品质测量链与硬测量链 \\
\bottomrule
\end{tabularx}
\end{table}

M12 若缺少两类食材的等温线参数区间，改为画 GAB 三个参数各自的单调方向，不画假数据点。M17 的冷点位置来自边界设定，不来自未运行的云图。

\section{数学图}

\begin{table}[H]
\centering
\caption{数学图 Q1--Q16}
\begin{tabularx}{\textwidth}{clX}
\toprule
编号 & 归属 & 图中必须出现的对象 \\
\midrule
Q1 & 命题 \ref{prop:steady} & 平衡点处灵敏度列消失 \\
Q2 & 命题 \ref{prop:pulse} & 两段脉冲后的两列灵敏度 \\
Q3 & 命题 \ref{prop:budget} & $\tau/\tau_2$ 与 $\Delta T_2$ 上界 \\
Q4 & 命题 \ref{prop:invalid} & 分位、阈值与滞回带 \\
Q5 & 命题 \ref{prop:shield} & 可行域不随后验变形 \\
Q6 & 命题 \ref{prop:path} & 链式法则的三支 \\
Q7 & 第 4 章 & $\mathrm{Fo}$、$\mathrm{Bi}$、$\mathrm{Ko}$ 的区域划分 \\
Q8 & 第 5 章 & 逆热传导误差放大随频率与厚度 \\
Q9 & 第 5 章 & 可观测矩阵秩随 $R_{\mathrm{int}}/R_{\mathrm{ext}}$ 恶化 \\
Q10 & 第 5 章 & 湿度饱和段的信息量下降 \\
Q11 & 第 6 章 & $G(a)-\lambda C(a)$ 在探针库上的比较结构 \\
Q12 & 第 7 章 & 干预与条件相关的差别：策略边被切断 \\
Q13 & 第 8 章 & 残差界 $r_{\max}$ 与回退 \\
Q14 & 第 9 章 & 机会约束与尾部硬约束的包含关系 \\
Q15 & 第 9 章 & 预测时域与控制时域 \\
Q16 & 第 10 章 & T 与 E 的数据隔离 \\
\bottomrule
\end{tabularx}
\end{table}

Q1--Q6 可以在证明完成时绘制，因为它们只含命题中的几何关系。Q8、Q9 需要数值秩，留到计算完成；图位先给出坐标。

\section{数据图}

\begin{table}[H]
\centering
\caption{数据图 D1--D24}
\begin{tabularx}{\textwidth}{clX}
\toprule
编号 & 坐标 & 必须包含的对照 \\
\midrule
D1 & 时间 / $T_1,T_2$ & 真值与分位带 \\
D2 & 时间 / $\varepsilon_l$ 表面与平均 & 蒸、烤、组合 \\
D3 & 时间 / $\alpha_m,\alpha_b$ & 同一温度史的分歧 \\
D4 & 时间 / $S,R_{\mathrm{int}}$ & 结壳反馈 \\
D5 & 空间 / 终止时温度 & 冷点位置 \\
D6 & 脉冲编号 / $\log\det\Sigma$ & P0、P1、P2、P4 \\
D7 & 厚度 / $G$ & 信息上限与停机需求 \\
D8 & 含水 $\times$ 厚度 / $G$ & 等高线 \\
D9 & 厚度 / 路径 A、B、C & 蒸汽干预 $+10\%$ \\
D10 & 策略 / $e_{\alpha}$ & 四策略分布 \\
D11 & 策略 / 能量与时间 & 次要指标，不单独下结论 \\
D12 & 策略 / $U_T$ & 均匀性 \\
D13 & 有无质量通道 / $e_T$ & 消融 \\
D14 & 有无热像 / $e_T$ & 消融 \\
D15 & 热像遮挡时长 / 探针有效率 & 蒸制故障 \\
D16 & 噪声倍数 / $e_{\alpha}$ & 稳健性 \\
D17 & 算子种类 / 干预位移误差 & FNO 与 DeepONet \\
D18 & 残差界 / 回退次数 & 有界更新 \\
D19 & 时间 / 请求加热与投影加热 & 贴近 $h_T=0$ \\
D20 & 是否隔离 / 越界次数 & 去掉屏障的消融 \\
D21 & 同构与异构 / $e_{\alpha}$ & 防止把自洽当成感知 \\
D22 & 厚度 50 mm / 区间宽度 & 拒绝单点剩余时间 \\
D23 & 种子 / 停止时刻 & 可复现散点 \\
D24 & 降级级别 / 停留时间 & 状态机是否被触发 \\
\bottomrule
\end{tabularx}
\end{table}

D10 是主图。D11 不得在 D10 缺席时单独出现。D21 中同构结果不进入摘要句。

\section{实施例图}

\begin{table}[H]
\centering
\caption{实施例图 E1--E8}
\begin{tabularx}{\textwidth}{clX}
\toprule
编号 & 内容 & 约束 \\
\midrule
E1 & 腔体截面：加热、蒸汽、风、排汽 & 只标模型用到的量 \\
E2 & 食品侧：质量、热像视场 & 穿刺探针画在对照侧 \\
E3 & 安全侧：漏电、水位、管温、门、腔压 & 与 E2 分开 \\
E4 & P4 在执行器上的时序 & 与 M18 一致 \\
E5 & 界面上的两项路径数值 & 与 D9 的定义一致 \\
E6 & 投影前后的加热指令 & 与 D19 一致 \\
E7 & 降级后面板上的区间显示 & 不出现单点分钟数 \\
E8 & 数据文件与页面四块的对应 & 列名与第 \ref{ch:demo} 章相同 \\
\bottomrule
\end{tabularx}
\end{table}

76 幅图的文字说明计入成稿篇幅。每幅说明包含：变量、对照、日志或证明来源、该图不能支持的句子。仅插图不加说明，不计入技术内容。

S1 与 S2 还要求下列图作为 D 图的扩展面板，编号 D25--D36。它们仍从同一日志取列，不新增科学主张。

\begin{table}[H]
\centering
\caption{扩展数据面板}
\begin{tabularx}{\textwidth}{llX}
\toprule
编号 & 坐标 & 用途 \\
\midrule
D25 & 原型 $\times$ 策略 / $e_{\alpha}$ 中位 & 排序是否依赖食材 \\
D26 & 非均质档 / 冷点离几何中心的距离 & T2 相对 T1 \\
D27 & 含水偏差 / $G_{\mathrm{real}}$ & 先验误差是否吞掉探针 \\
D28 & 噪声倍数 / 探针有效率 & S2 \\
D29 & 遮挡时长 / 区间宽度 & 拒绝单点时间的频率 \\
D30 & 脂肪层开关 / 路径 C & 慢路径是否出现 \\
D31 & $E_a$ 扫描 / $\alpha_m$ & 动力学数量级的影响 \\
D32 & 集合 32 与 128 / 方差比 & 在线集合是否过自信 \\
D33 & T3 风速 / $h$ & 边界库的形状 \\
D34 & 种子分位 / 停止时刻差 B3$-$B2 & 探针贡献的分布 \\
D35 & 扩充标记 / 干预误差 & 防止教师扩充进入主表 \\
D36 & 来源字段 / 参数轴 & 区分 DOI、文献区间、仿真 \\
\bottomrule
\end{tabularx}
\end{table}

至此图位为机制 28、数学 16、数据 36、实施例 8，共 88。D25--D36 未有日志时同样只留坐标。
